\begin{theorem}[MetaCIC parallel-diamond frontier seed-status public export]
\label{thm:metacic-parallel-diamond-frontier-seed-status-public-export}
The paper-side seed-status reader for $\MetaCICParallelDiamondFrontierUp$ is consumed through the current route of \autoref{ch:concrete-instances-metacic-parallel-diamond-frontier-namecert} and the binding status block in \texttt{closurestatus\_current.tex}. A public reader may cite that route only as the displayed seed-status export over the finite packet
$$
P,K,J,R,S,C,B,O,H,T,G,N .
$$
It is not a residual-substitution discharge, unrestricted subject-reduction theorem, Church--Rosser theorem, evaluator equality theorem, quotient equality theorem, global normalization theorem, or kernel-consistency theorem.
\end{theorem}
\begin{proof}
Read the hub through the existing seed-status children and then through the current closurestatus child. These files expose the retained residual-substitution premise, critical-pair row, candidate-join row, residual-frontier row, closed-checker row, closed-normal fragment, bounded-check row, obstruction row, transport, replay, provenance, and local naming before any status read is available.

The closurestatus block supplies the manuscript's current public status route for that same packet. Since the route has no carrier coordinate for residual-substitution discharge, unrestricted subject reduction, Church--Rosser, evaluator equality, quotient equality, global normalization, or kernel consistency, the public export cannot be reread as any of those endpoints.
\end{proof}
